\section*{\S~2c Stab und erweiterter Stab}
\label{sec:stab_und_erweiterter_stab}
\addcontentsline{toc}{section}{\S~2c Stab und erweiterter Stab}

\begin{enumerate}[label=(\arabic*)]
    \item Der Stab und der erweiterte Stab der 1st Aircat Division unterstützen die Leitung in allen administrativen Aufgaben und Entscheidungsprozessen.
    \item Der Stab tagt bei Bedarf, mindestens aber einmal im Quartal. Der erweiterte Stab tagt ebenfalls bei Bedarf, mindestens aber einmal im Halbjahr.
    \item Der Stab besteht aus der Leitung, einem Gefreiten vom Dienst (GvD) in einem Manschaftsdienstgrad ab Gefreiter und, sofern mindestens ein Mitglied den Dienstgrad Unteroffizier oder Stabsunteroffizier innehat, einem Unteroffizier vom Dienst (UvD) im Rang des Unteroffiziers oder Stabsunteroffiziers. Der erweiterte Stab besteht aus dem Stab sowie allen weiteren aktiven Unteroffizieren mit Portepee und Unteroffizieren ohne Portepee. Die Mitglieder des Stabes verzichten, sofern nicht anders festgelegt, während ihrer Dienstzeit auf ihr Recht, Einspruch gegen Entscheidungen des Stabes oder des erweiterten Stabes einzulegen. Die Mitglieder des erweiterten Stabes verzichten, sofern nicht anders festgelegt, während ihrer Dienstzeit auf ihr Recht, Einspruch gegen Entscheidungen des erweiterten Stabes einzulegen.
\item Der Gefreite vom Dienst (GvD) wird immer in der Nachbesprechung der letzten Mission in einem Quartal oder im Rahmen einer Mitgliederversammlung durch die aktive Mannschaft gewählt und tritt seinen Dienst zum ersten Tag des neuen Quartals an.

\begin{enumerate}[label=(\Alph*)]

    \item Der Gefreite vom Dienst (GvD) repr\"{a}sentiert w\"{a}hrend seiner Amtszeit die Mannschaft. Er verzichtet w\"{a}hrend seiner Amtszeit auf sein Recht, Beschwerde gegen Einzelpersonen, die einen Mannschaftsdienstgrad innehaben, einzureichen oder Einspruch gegen Entscheidungen des Stabes einzulegen, sofern ihm dies nicht durch ein Mitglied in einem Mannschaftsdienstgrad aufgetragen wurde oder es sich um einen schweren Regelversto\ss{} nach \S~8 Absatz \ref{itm:schwerer_verstoß} handelt.
    
    \item Die Unit-Leitung hat f\"{u}r diese Wahl den Posten der Wahleitung inne.
    
    \item Die Wahlleitung erinnert die Mitglieder der Unit zu einem geeigneten Zeitpunkt, sp\"{a}testens jedoch in der Nachbesprechung der letzten Mission vor der Wahl an die bevorstehende Wahl des Gefreiten vom Dienst (GvD).

    \item Die Wahl in der Nachbesprechung der letzten Mission in einem Quartal kann auch in einem hybriden Verfahren durchgef\"{u}hrt werden. In diesem Verfahren werden die Vorschl\"{a}ge in Anwesenheit abgegeben, w\"{a}hrend die eigentliche Wahl per 72-st\"{u}ndiger Abstimmung \"{u}ber eine Drittplattform durchgef\"{u}hrt wird.

    \item Eine Wahl in Abwesenheit ist zul\"{a}ssig. Vorschl\"{a}ge f\"{u}r Mitglieder, die nicht anwesend sind, m\"{u}ssen sp\"{a}testens zu Beginn der Mission, in deren Nachbesprechung die Wahl stattfindet, oder zu Beginn der Mitgliederversammlung bei der Wahlleitung eingegangen sein.

    \item Mitglieder, die den Dienstgrad Korporal oder Stabskorporal innehaben, k\"{o}nnen nur durch eigenen Vorschlag und unter Ber\"{u}cksichtigung ihrer sonstigen Verpflichtungen, die sie durch ihren Dienstgrad innehaben, f\"{u}r den Posten des Gefreiten vom Dienst (GvD) kandidieren.

    \item Jede vorgeschlagene Person hat das Recht, ihrer Kandidatur zu widersprechen. In diesem Fall ist die Kandidatur unwirksam.

    \item Es gibt keine Obergrenze f\"{u}r Amtszeiten, allerdings d\"{u}rfen keine zwei Amtszeiten aufeinander folgen.

    \item Der bisherige Gefreite vom Dienst (GvD) bleibt so lange in seinem Amt, bis ein neuer Gefreiter vom Dienst (GvD) gew\"{a}hlt wurde. 
    
    \end{enumerate}    
\item Der Unteroffizier vom Dienst (UvD) wird immer in der Nachbesprechung der letzten Mission in einem Quartal oder im Rahmen einer Mitgliederversammlung durch die Leitung ernannt und tritt seinen Dienst zum ersten Tag des neuen Quartals an. Es gibt keine Obergrenze für Amtszeiten, allerdings dürfen keine zwei Amtszeiten aufeinander folgen, solange sich dies durch die Anzahl an Unteroffizieren vermeiden lässt.
\end{enumerate}